\begin{thebibliography}{99}

\bibitem{nash1994}
W. J. Nash, T. L. Sellers, S. R. Talbot, A. J. Cawthorn, and W. B. Ford, ``The population biology of abalone (\emph{Haliotis} species) in Tasmania. I. Blacklip abalone (\emph{H. rubra}) from the north coast and islands of Bass Strait,'' Sea Fisheries Division, Tasmania, Tech. Rep. 48, 1994.

\bibitem{uci}
W. Nash, T. Sellers, S. Talbot, A. Cawthorn, and W. Ford, ``Abalone,'' UCI Machine Learning Repository, 1994. [Online]. Available: https://doi.org/10.24432/C55C7W

\bibitem{waugh1995}
S. G. Waugh, ``Extending and benchmarking Cascade-Correlation: Extensions to the Cascade-Correlation architecture and benchmarking of feed-forward supervised artificial neural networks,'' Ph.D. dissertation, Dept. Comput. Sci., Univ. Tasmania, Hobart, Australia, 1995.

\bibitem{cybenko1989}
G. Cybenko, ``Approximation by superpositions of a sigmoidal function,'' \emph{Math. Control, Signals, Syst.}, vol. 2, no. 4, pp. 303--314, 1989.

\bibitem{hornik1989}
K. Hornik, M. Stinchcombe, and H. White, ``Multilayer feedforward networks are universal approximators,'' \emph{Neural Netw.}, vol. 2, no. 5, pp. 359--366, 1989.

\bibitem{rumelhart1986}
D. E. Rumelhart, G. E. Hinton, and R. J. Williams, ``Learning representations by back-propagating errors,'' \emph{Nature}, vol. 323, no. 6088, pp. 533--536, 1986.

\bibitem{bengio2012}
Y. Bengio, ``Practical recommendations for gradient-based training of deep architectures,'' in \emph{Neural Networks: Tricks of the Trade}, 2nd ed., G. Montavon, G. B. Orr, and K.-R. M\"{u}ller, Eds. Berlin, Germany: Springer, 2012, pp. 437--478.

\bibitem{goodfellow2016}
I. Goodfellow, Y. Bengio, and A. Courville, \emph{Deep Learning}. Cambridge, MA, USA: MIT Press, 2016.

\bibitem{bottou2010}
L. Bottou, ``Large-scale machine learning with stochastic gradient descent,'' in \emph{Proc. COMPSTAT'2010}, Paris, France, 2010, pp. 177--186.

\bibitem{kingma2015}
D. P. Kingma and J. Ba, ``Adam: A method for stochastic optimization,'' in \emph{Proc. 3rd Int. Conf. Learn. Represent. (ICLR)}, San Diego, CA, USA, 2015.

\bibitem{wilson2017}
A. C. Wilson, R. Roelofs, M. Stern, N. Srebro, and B. Recht, ``The marginal value of adaptive gradient methods in machine learning,'' in \emph{Advances in Neural Information Processing Systems 30}, 2017, pp. 4148--4158.

\bibitem{tukey1977}
J. W. Tukey, \emph{Exploratory Data Analysis}. Reading, MA, USA: Addison-Wesley, 1977.

\bibitem{lecun1998}
Y. LeCun, L. Bottou, G. B. Orr, and K.-R. M\"{u}ller, ``Efficient BackProp,'' in \emph{Neural Networks: Tricks of the Trade}, G. B. Orr and K.-R. M\"{u}ller, Eds. Berlin, Germany: Springer, 1998, pp. 9--50.

\bibitem{pedregosa2011}
F. Pedregosa \emph{et al.}, ``Scikit-learn: Machine learning in Python,'' \emph{J. Mach. Learn. Res.}, vol. 12, pp. 2825--2830, 2011.

\bibitem{nair2010}
V. Nair and G. E. Hinton, ``Rectified linear units improve restricted Boltzmann machines,'' in \emph{Proc. 27th Int. Conf. Mach. Learn. (ICML)}, Haifa, Israel, 2010, pp. 807--814.

\bibitem{glorot2010}
X. Glorot and Y. Bengio, ``Understanding the difficulty of training deep feedforward neural networks,'' in \emph{Proc. 13th Int. Conf. Artif. Intell. Statist. (AISTATS)}, Sardinia, Italy, 2010, pp. 249--256.

\bibitem{he2009}
H. He and E. A. Garcia, ``Learning from imbalanced data,'' \emph{IEEE Trans. Knowl. Data Eng.}, vol. 21, no. 9, pp. 1263--1284, 2009.

\bibitem{fawcett2006}
T. Fawcett, ``An introduction to ROC analysis,'' \emph{Pattern Recognit. Lett.}, vol. 27, no. 8, pp. 861--874, 2006.

\bibitem{hand2001}
D. J. Hand and R. J. Till, ``A simple generalisation of the area under the ROC curve for multiple class classification problems,'' \emph{Mach. Learn.}, vol. 45, no. 2, pp. 171--186, 2001.

\end{thebibliography}
